\documentclass[10pt,a4paper]{amsart}
\usepackage[margin=19mm]{geometry}
\usepackage{amsmath,amssymb,mathtools}
\usepackage[scr=rsfs,cal=euler]{mathalfa}
\usepackage{xfrac}
\usepackage[unicode,hidelinks,bookmarks=false]{hyperref}
\newcommand{\ie}{i.e.\ }
\newcommand{\cf}{cf.\ }
\newcommand{\resp}{respectively\ }
\newcommand{\Subsection}{\subsection}
\providecommand{\nobreakdash}{\nobreak\hbox{-}}
\setlength{\emergencystretch}{2em}
\multlinegap0pt
\usepackage{bm}
\usepackage{bbm}
\usepackage{dsfont}
\usepackage{xspace}
\usepackage{cancel}
\usepackage{mathtools}
\usepackage{amsthm}
\makeatletter
\@ifundefined{tabular}{\newtheorem{tabular}{Tabular}}{}
\@ifundefined{thm}{\newtheorem{thm}{Theorem}}{}
\makeatother
\newcommand{\OO}{\mathcal{O}}
\newcommand{\p}{\mathfrak{p}}
\title[On the fourth power level of $\mathfrak{p}$-adic completions]{On the fourth power level of $\mathfrak{p}$-adic completions of biquadratic number fields}
\author{Kazimierz Chomicz}
\date{Source: arXiv:2503.21559v3 (CC BY 4.0)}
\begin{document}
\maketitle

\noindent\emph{Source and licence.} On the fourth power level of $\mathfrak{p}$-adic completions of biquadratic number fields. Version arXiv:2503.21559v3 (CC BY 4.0), distributed under the Creative Commons Attribution 4.0 International licence, as recorded in the arXiv licence metadata for this exact version and shown on the abstract page https://arxiv.org/abs/2503.21559v3. Full rights notice: licenses/arxiv-2503.21559-CC-BY-4.0.txt.\par\smallskip

\noindent\textit{Editorial scope.} Counted unit: the Thm whose source label is \texttt{main2} in the article (source0.tex lines 843--871, source label \texttt{main2}), delivered together with the article's own complete original proof of it (source0.tex lines 872--1015, one \texttt{proof} environment of 144 lines). No mathematics is written, repaired or re-derived: statement and proof are the article's own text line for line. Required context retained uncounted: 1.1 (lines 130--132); 1.2 (lines 135--138); 1.3 (lines 142--148); main (lines 479--481). Every definition, display and label of those lines is retained unchanged. Omitted from this package and not redistributed: 1--129, 133--134, 139--141, 149--478, 482--842, 1016--1178. Nothing inside a retained excerpt is omitted or altered: every retained body is delivered exactly as the article's lines are written, so no omission receipt of any kind accompanies this package. Citations: the retained excerpts carry no citation command at all, so no bibliography entry of the article is transcribed and no reference is redistributed. Reference adaptation: none was needed and none was made; every reference inside the delivered unit resolves inside it, so no unresolved reference or citation is produced. Printed slips kept literal: no printed word, symbol, hypothesis or number of the counted statement or of its proof was altered. Layout and class: the article's own class is \texttt{article} with packages including inputenc, geometry, fancyhdr, lastpage, amsmath, amssymb, graphicx, enumitem, multicol, amsfonts, amsthm, tikz-cd; the wrapper is the lane's pinned \texttt{amsart} editorial wrapper at 10pt on A4, which restates the theorem-like environments the retained text needs and adds the article's own macro definitions that the retained text needs (\texttt{\textbackslash OO}, \texttt{\textbackslash p}), with extra wrapper packages (bm, bbm, dsfont, xspace, cancel, mathtools). The delivered numbering is the unit's own. Assets: no retained span contains a figure environment, a graphics-inclusion command, a PDF-page inclusion, a table or a TikZ picture; no image file is copied into this package, the package declares no asset dependency and no image file is redistributed with it. Licence: version arXiv:2503.21559v3 of this article is distributed under the Creative Commons Attribution 4.0 International licence, as recorded in the arXiv licence metadata for this exact version and shown on the abstract page https://arxiv.org/abs/2503.21559v3. A full rights notice is delivered at licenses/arxiv-2503.21559-CC-BY-4.0.txt. Characters: every retained excerpt is pure ASCII, so the delivered engine, XeLaTeX with its default Latin Modern text font, reports no missing character.
\begin{thm}\label{1.1} 
We have an equivalence  $e=f=1 \iff s=15$. 
\end{thm}

\begin{thm}\label{1.2} 
 If $f$ is even then $s \leq 2$. If $s=1$ then $4 \mid e$. 

\end{thm}

\begin{thm}\label{1.3} 
 Let $e=2,\ f=1$. Let $\pi$ be any element of $\mathcal{O}_K $ such that $\mathfrak{p} \parallel \pi$ $($i.e. $\pi \in \mathfrak{p} \setminus \mathfrak{p}^2$ $)$. \\
Then $$s = \begin{cases}
    6 \text{ if } \pi^2 \equiv 2 \pmod{\mathfrak{p}^4}\\
    4 \text{ otherwise.}
\end{cases} $$
\end{thm}

\begin{thm}\label{main}
Let $K$ be a number field, with a prime ideal $\p \mid (2)$ of $\OO_k$ with $e(\p) = 4, \ f(\p)=1$. Let $\tau$ be any element in $\mathfrak{p}\setminus \mathfrak{p}^2$. Then $s = s_4(K_{\mathfrak{p}})$ is one of $\{1,2,3,4,5,6\}$, specifically
\end{thm}

\begin{thm}\label{main2} Let $m,n$ be two square free integers. Let $K = \mathbb{Q}(\sqrt{m},\sqrt{n})$, $\mathfrak{p} \mid (2)$, $d = \gcd(m,n)$, $k = \frac{mn}{d^2}$. Then $s = s_4(K_{\mathfrak{p}})$ is one of $\{1,2,3,4,6,15\}$, specifically $($up to rearrangement of $m,n,k)$
\begin{center}
    if $m \equiv 3$, $n\equiv k \equiv 2 \pmod{4}$ then $s = \begin{cases}
        \textbf{1} \text{ if } 32 \mid n+k \text{ and } \frac{nk}{4} \equiv 15 \pmod{16} \\
        \text{or } 16 \parallel n+k \text{ and } \frac{nk}{4} \equiv 7 \pmod{16}\\
        \textbf{2} \text{ if } 32 \mid n+k \text{ and } \frac{nk}{4} \equiv 7 \pmod{16} \\
        \text{or } 16 \parallel n+k \text{ and } \frac{nk}{4} \equiv 15 \pmod{16}\\
        \textbf{3} \text{ if } 8 \parallel n+k
    \end{cases}$
\end{center}
\begin{center}
    if $m \equiv 1$, $n \equiv k \equiv 3 \pmod{4}$ then $s = \begin{cases}
        \textbf{4} \text{ if } m \equiv 1 \pmod{8}\\
        \textbf{2} \text{ if } m \equiv 5 \pmod{8}
    \end{cases}$
\end{center}
\begin{center}
    if $m \equiv 1$, $n \equiv k \equiv 2 \pmod{4}$ then $s = \begin{cases}
        \textbf{6} \text{ if } m \equiv 1 \pmod{8}\\
        \textbf{2} \text{ if } m \equiv 5 \pmod{8}
    \end{cases}$
\end{center}
\begin{center}
    if $m \equiv n \equiv k \equiv 1 \pmod{4}$ then $s = \begin{cases}
        \textbf{15} \text{ if } m \equiv n \equiv k \equiv 1 \pmod{8}\\
        \textbf{2} \text{ otherwise.}
    \end{cases}$
\end{center}
\end{thm}

\begin{proof}
Again, proof will be a case by case analysis. Moreover some computations were conducted in Macaulay2, especially when we need to explicitly know the structure of finite rings and its prime ideals.



\textbf{Case 3.1.} $m \equiv 3$, $n \equiv k \equiv 2 \pmod{4}$ and the basis of $\mathcal{O}_K$ is $\{1, \sqrt{m}, \sqrt{n}, \frac{\sqrt{n}+\sqrt{k}}{2}\} = \{1,a,b,c\}$. \\
We start with representing the following elements as integral linear combinations of the basis:
\begin{center}
\begin{tabular}{@{}l@{}}
    $a^2 = m \equiv 1 \pmod{2}$\\
    $b^2 = n \equiv 0 \pmod{2}$\\
    $c^2 = \frac{n+k}{4} + \frac{n}{2d}\sqrt{m} \equiv a \pmod{2}$\\
    $ab = d\sqrt{k} = 2d \cdot \frac{\sqrt{n} + \sqrt{k}}{2} - d\sqrt{n} \equiv b \pmod{2}$\\
    $ac = \frac{d\sqrt{k}+ \frac{m}{d} \sqrt{n}}{2} = d \cdot \frac{\sqrt{n} + \sqrt{k}}{2} + \frac{\frac{m}{d} - d}{2} \sqrt{n} \equiv c + b \pmod{2}$\\
    $bc = \frac{n + \frac{n}{d} \sqrt{m}}{2} \equiv 1 + a \pmod{2}.$
\end{tabular}
\end{center}
In the above calculations, we have used the fact that $8 \mid n+k$. If this is not the case then $n \equiv k \pmod{8}$, but $n \equiv k \equiv \frac{mn}{d^2} \pmod{8} \implies d^2 \equiv m \equiv 3 \pmod{4} $ a contradiction. 

We have obtained the ideal $I = (x^2+1,y^2,z^2+x,xy+y,xz+z+y,yz+1+x)$ with already reduced coefficients modulo 2. constructed from the above remainders modulo 2, we have that $\mathbb{F}_2[x,y,z]/I \cong \mathcal{O}_K/(2)$. Using the calculator Macaulay2 we will find prime ideals $P_1, \dots , P_r$ of $\mathbb{F}_2[x,y,z]/I$ and check what is the smallest power $\alpha$, so that $(P_1 \cdots P_r)^{\alpha} =(0)$. This will imply that $I = (P_1 \cdots P_r)^{\alpha}$ and $(2) = (\mathfrak{p}_1 \cdots \mathfrak{p}_r)^{\alpha}$ (the ramification indices must be equal, because biqadratic extension is Galois). By the introduction, we know what are the possible factorizations.


We will now proceed with computations in Macaulay2. At the very end we present first case on how we did it. Hence, in order to shorten the paper we will avoid repeating the code in each case.

For the ideal $I$ we get that the ideal $J =  (x+1,y,z+1)$ is prime with $J^2 = (x+y+1,x+1)$, $J^3 = (x+y+1)$ and $J^4= 0$ all in $\mathbb{F}_2[x,y,z]/I$. Here, $x,y,z$ of course denote their corresponding images in the quotient ring. $J^4= 0$ implies $(2) = \mathfrak{p}^4$, hence $e=4,f=1$ in $K$. We have  $z+1 \in J$ and $z+1 \not\in J^2$ hence, as the generator of $\p$ we may take $ c+1 = \tau \in \mathfrak{p}\setminus\mathfrak{p}^2$. Now, using $\tau = c+1$, $c^2 \equiv a$ and $a^2 = m \equiv 3 \pmod{4}$ we have to check which case from Theorem \ref{main} we can apply. See that $\tau^4 = (c+1)^4 = c^4+4c^3+6c^2+4c+1 \equiv c^4+2c^2+1 \equiv (c^2+1)^2 \equiv  (a+1)^2 \equiv a^2+2a+1 \equiv 2a \pmod{4 = \mathfrak{p}^8}$. Thus $\tau^4+2 \equiv 2(a+1) \pmod{4 = \mathfrak{p}^8}$. Now $x+1 \in J^2$ and $x+1 \not\in J^3$ implies $\mathfrak{p}^2 \parallel a+1 \implies \mathfrak{p}^6 \parallel \tau^4+2$.


By Theorem \ref{main} the fourth level $s$ depends on remainders modulo $\mathfrak{p}^{13}$ of the following:
\begin{center}
\begin{tabular}{@{}l@{}}
    $A_1 = 4\tau^3+4\tau^2+4\tau+(\tau^4+2\tau^2+2),$\\
    $A_2 = 4\tau + (\tau^4+1)(\tau^4+2\tau^2+2),$\\
    $A_3 = 4\tau + (\tau^4+2\tau^2+2), $\\
    $A_4= 4\tau^3+4\tau^2+4\tau+(\tau^4+1)(\tau^4+2\tau^2+2). $
\end{tabular}
\end{center}
Computation of $A_2,A_3,A_4$ will be based on the value of $A_1$.
\[\tau^4+2\tau^2+2 = (c+1)^4 + 2(c+1)^2+2 = c^4 + 4c^3 + 8c^2+8c+5 \equiv c^4+ 4c^3+8a+8c+5 \pmod{\mathfrak{p}^{13}},\]
\[4\tau^3 + 4\tau^2 + 4\tau = 4c^3+16c^2+24c+12 \equiv 4c^3+8c+12 \pmod{\mathfrak{p}^{13}} \]
\[ \implies A_1 = 4\tau^3+4\tau^2+4\tau+(\tau^4+2\tau^2+2) \equiv c^4+8c^3+8a+1 \equiv   \frac{(n+k)^2}{16} + \frac{nk}{4} + \frac{\frac{n}{d}(n+k)a}{4} + 8(c+b) + 8a + 1 \]
\[ \equiv \frac{(n+k)^2}{16} + \frac{nk}{4} + \frac{\frac{n}{d}(n+k)a}{4} + 1  \pmod{\mathfrak{p}^{13}}. \]
In the last congruence we used the fact that $8(a+b+c) \equiv 0 \pmod{\mathfrak{p}^{13}}$, and this follows from $(a+b+c) \equiv (a+b+c) + (a+b+1) \equiv c+1 \equiv 0 \pmod{\mathfrak{p}}$. We now have 2 cases:
\\ \\
$1^{\circ}$ $16 \mid n+k$. In this case $\frac{nk}{4} \equiv 7 \pmod{8}$ and $A_1 \equiv \frac{nk}{4} + \frac{\frac{n}{d}(n+k)a}{4} + 1  \pmod{\mathfrak{p}^{13}}$, hence $A_1 \equiv 7 + \frac{32a}{4}+ 1 \equiv 0 \pmod{8 = \mathfrak{p}^{12}}$. This readily implies $\mathfrak{p}^{9} \parallel \tau^4+2\tau^2+2$. This on the other hand implies $A_1 \equiv A_4, \  A_2 \equiv A_3 \pmod{\mathfrak{p}^{13}}$. We have $A_2 \equiv A_1 + 4\tau^2+4\tau^3 \pmod{\mathfrak{p}^{13}}$ which implies that $\mathfrak{p}^{10} \parallel A_2$. This gives that the value of $s$ depends only on $A_1 \pmod{\mathfrak{p}^{13}}$ - if $A_1 \equiv 0$ then $s=1$ and if not, then $\mathfrak{p}^{12} \parallel A_1$, so $s=2$ in this case.



If $32 \mid n+k$ and $\frac{nk}{4} \equiv 15 \pmod{16}$ or $16 \parallel n+k$ and $\frac{nk}{4} \equiv 7 \pmod{16}$ we have $A_1 \equiv 15 + 1 \equiv 0 \pmod{\p^{13}}$ or $A_1 \equiv 7 + 8a + 1 \equiv 8(a+1) \pmod{\p^{13}}$, respectively. In both cases we get $s=1$. 

If $32 \mid n+k$ and $\frac{nk}{4} \equiv 7 \pmod{16}$ or $16 \parallel n+k$ and $\frac{nk}{4} \equiv 15 \pmod{16}$ we have $\mathfrak{p}^{12} \parallel A_1$, hence $s=2$.
\\ \\
$2^{\circ}$ $8 \parallel n+k$. In this situation we have $\frac{nk}{4} \equiv 3 \pmod{8}$. Also, $A_1 \equiv 4 + 3 + 4a + 1 \equiv 4a \pmod{8 = \mathfrak{p}^{12}}$. Since $a \not\in \mathfrak{p}$ we have $\mathfrak{p}^8 \parallel A_1$. This implies that $\mathfrak{p}^8 \parallel A_i$ for $i=2,3,4$. By Theorem \ref{main}, to determine $s$ we need to check if any of $\{A_1+\tau^8, A_2+\tau^8, A_3+\tau^8, A_4+\tau^8\}$ is congruent to $0 \pmod{\mathfrak{p}^{12} = 8}$. We have the following chain of congruences:
\[ \tau^8 = (c+1)^8 = (c^2+2c+1)^4 = (c^2+1  + 2\gamma)^4  \equiv (a+1)^4 \]
\[  \equiv a^4 + 4a^3+6a^2+4a+1 \equiv m^2 + 4a(m+1) + 6m + 1 \equiv 1 + 18 + 1 \equiv 4 \pmod{8}. \]
Furthermore, since $a+1, \  a+b+1 \not\in (2)$ we obtain
\[ A_1 + \tau^8 \equiv 4a + 4 \equiv 4(a+1) \not\equiv 0 \pmod{8 },\]
\[A_3 + \tau^8 \equiv A_1 + \tau^8 + 4\tau^2+4\tau^3 \equiv 4a + 4 + 4(c+1)^2 + 4(c+1)^3 \equiv 4a + 4c^3 + 16c^2 + 20c + 12\]  \[\equiv 4a + 4(c+b) + 4c + 4 \equiv 4(a+b+1) \not\equiv 0 \pmod{8 }.\]
Also $A_1 \equiv A_4$ and $A_2 \equiv A_3 \pmod{8 = \mathfrak{p}^{12} }$, because $\mathfrak{p}^8 \parallel \tau^4+2\tau^2+2$ . Hence $s=3$ by Theorem \ref{main}.
\\ \\
\textbf{Case 3.2.} $m \equiv 1$, $n \equiv k \equiv 3 \pmod{4}$ and the basis of $\mathcal{O}_K$ is $\{1, \frac{1+\sqrt{m}}{2}, \sqrt{n}, \frac{\sqrt{n}+\sqrt{k}}{2}\} = \{1,a,b,c\}$. As before we calculate:
\begin{center}
\begin{tabular}{@{}l@{}}
$a^2 = \frac{m-1}{4} + \frac{1+\sqrt{m}}{2} \equiv a \text{ or } 1 + a \pmod{2}$ \\
$b^2 = n \equiv 1 \pmod{2}$\\
$c^2 = \frac{n+k+2\sqrt{nk}}{4} = \frac{n+k}{4} + \frac{n}{2d}\sqrt{m} = \frac{n+k-2\frac{n}{d}}{4} + \frac{n}{d} \cdot \frac{1+\sqrt{m}}{2} \equiv a \text{ or } 1+a \pmod{2}$\\
$ab = \frac{\sqrt{n} + d\sqrt{k}}{2} = d \cdot \frac{\sqrt{n}+\sqrt{k}}{2} - \frac{d-1}{2} \sqrt{n} \equiv c \text{ or } c + b \pmod{2} $\\
$ac = \frac{\sqrt{n} + \sqrt{k} + d\sqrt{k} + \frac{m}{d} \sqrt{n} }{4} =  \frac{d+1}{2} \cdot \frac{\sqrt{n}+\sqrt{k}}{2} + (\frac{\frac{m}{d}-d}{4})\sqrt{n} \equiv b+c \text{ or } c \text{ or } b \text{ or } 0 \pmod{2}$\\
$bc = \frac{n + \frac{n}{d}\sqrt{m}}{2} = \frac{n - \frac{n}{d}}{2} + \frac{n}{d} \cdot \frac{1+ \sqrt{m}}{2} \equiv a \text{ or }  1 + a \pmod{2}.$
\end{tabular}
\end{center}
The 'or' in $a^2$ depends on $m \pmod{8}$. The 'or' in $c^2$ depends on $n+k-2\frac{n}{d} \pmod{8}$. The 'or'-s in $ac$ depend on $m \pmod{8}$ and $d \pmod{4}$ and the remaining cases depend on $d \pmod{4}$. Note that $k \equiv mn \pmod{8}$, so by $d \pmod{4}$ and $m \pmod{8}$ we can compute $n+k-2\frac{n}{d} \equiv n+k-2\frac{3}{d} \pmod{8}$. If $m\equiv 1\pmod{8}$, then $n \equiv k \equiv 3 \pmod{4}$ and $k \equiv mn \equiv n \pmod{8}$ gives $n+k \equiv 6 \pmod{8}$. If $m \equiv 5 \pmod{8}$ then $n+k \equiv 2 \pmod{8}$. Hence we have to check 4 cases.


(a) $m \equiv 1 \pmod{8}$, $d \equiv 1 \pmod{4}$ $ \implies n+k-2\frac{n}{d} \equiv 0 \pmod{8}$. Similarly as before, we obtain an ideal $I = (x^2+x,y^2+1,z^2+x,xy+z,xz+z,yz+x)$. Again, using Macaulay2 we get the following data: ideals $A = (z,y+1,x)$ and $B = (z+1,y+1,x+1)$ are prime and $A^2 = (x,x+z,z)$, $B^2 = (x+1,x+y+z+1)$, $AB = (x+z,x+y+z+1)$ and $(AB)^2 = 0$. This all implies that $(2)=(\mathfrak{p}_1\mathfrak{p}_2)^2$ hence $ e=2, \ f=1$. Also, note that $y+1 \in A$ and $y+1 \not\in A^2$, hence $b+1 = \pi \in \mathfrak{p}_1 \setminus \mathfrak{p}_1^2$ and by Theorem \ref{1.3} $s = s_4(K_{\mathfrak{p}_1})$ depends on $\pi^2 = (b+1)^2 \pmod{\mathfrak{p}_1^4}$. We have $n \equiv 3 \pmod{\mathfrak{p}_1^4}$ and $y \not\in A \implies b = \sqrt{n} \not\in \mathfrak{p}_1$. 

Moreover, we have a chain of equalities and congruences $(b+1)^2 = (\sqrt{n}+1)^2 = n+1+2\sqrt{n} \equiv 2\sqrt{n} \not\equiv 2 \pmod{\mathfrak{p}_1^4}$. To justify this, see that there is an equivalence $2\sqrt{n} = 2b \equiv 2 \iff 2(b+1) \equiv 0 \pmod{\mathfrak{p}_1^4}$ which is does not hold since $\mathfrak{p}_1 \parallel b+1$, so by Theorem \ref{1.3} we have $s=4$ for $K_{\mathfrak{p}_1}$. For $\mathfrak{p}_2$ the situation is the same since the extension $K/\mathbb{Q}$ is Galois.
\\ \\
(b) $m \equiv 1 \pmod{8}$, $d \equiv 3 \pmod{4}$ $ \implies n+k-2\frac{n}{d} \equiv 4 \pmod{8}$. In this case we start with
$I = (x^2+x,y^2+1,z^2+x+1,xy+z+y,xz,yz+x+1)$ and we obtain prime ideals $A=(z,y+1,x+1)$ and $B = (z+1,y+1,x)$ satisfying $A^2 = (x+1,x+z+1,z)$, $B^2=(x,x+y+z)$, $AB = (x+z+1,x+y+z)$ and $(AB)^2=0$. As before we have $(2) = (\mathfrak{p}_1\mathfrak{p}_2)^2 \implies e=2, \ f=1$ and $b+1 = \pi \in \mathfrak{p}_1\setminus\mathfrak{p}_1^2$ same reasoning as in (a) yields $s=4$.
\\ \\
(c) $m \equiv 5 \pmod{8}$, $d \equiv 1 \pmod{4} \implies n+k-2\frac{n}{d} \equiv 4 \pmod{8}$. This time we have an ideal
$I = (x^2+x+1,y^2+1,z^2+x+1,xy+z,xz+z+y,yz+x)$ and a a single prime ideal $J = (y+1,x+z)$ satisfying $J^2 = 0.$ Hence, $ (2) = \mathfrak{p}^2 \implies e=2, \ f=2 \implies s=2$ by Theorem \ref{1.2}.
\\ \\
(d) $m \equiv 5 \pmod{8}$, $d \equiv 3 \pmod{4}$ $ \implies n+k-2\frac{n}{d} \equiv 0 \pmod{8}$. We get
$I = (x^2+x+1,y^2+1,z^2+x,xy+z+y,xz+y,yz+x+1)$ and a prime ideal $J = (y+1,x+z+1)$ satisfying $J^2 = 0.$ Thus $(2) = \mathfrak{p}^2 \implies e=2,\ f=2 \implies s=2$ by Theorem \ref{1.2}.
\\ \\
\textbf{Case 3.3.}  $m \equiv 1$, $n \equiv k \equiv 2 \pmod{4}$ and the basis of $\mathcal{O}_K$ is given by $\{1, \frac{1+\sqrt{m}}{2}, \sqrt{n}, \frac{\sqrt{n}+\sqrt{k}}{2}\} = \{1,a,b,c\}$.  As before, we calculate:
\begin{center}
\begin{tabular}{@{}l@{}}
$a^2 = \frac{m-1}{4} + \frac{1+\sqrt{m}}{2} \equiv a \text{ or } 1 + a\pmod{2}$, \\
$b^2 = n \equiv 0 \pmod{2}$,\\
$c^2 = \frac{n+k+2\sqrt{nk}}{4} = \frac{n+k}{4} + \frac{n}{2d}\sqrt{m} = \frac{n+k-2\frac{n}{d}}{4} + \frac{n}{d} \cdot \frac{1+\sqrt{m}}{2} \equiv 0 \pmod{2}$,\\

$ab = \frac{\sqrt{n} + d\sqrt{k}}{2} = d \cdot \frac{\sqrt{n}+\sqrt{k}}{2} - \frac{d-1}{2} \sqrt{n} \equiv c \text{ or } c+b \pmod{2}$,\\
$ac = \frac{\sqrt{n} + \sqrt{k} + d\sqrt{k} + \frac{m}{d} \sqrt{n} }{4} =  \frac{d+1}{2} \cdot \frac{\sqrt{n}+\sqrt{k}}{2} + (\frac{\frac{m}{d}-d}{4})\sqrt{n} \equiv 0 \text{ or } c \text{ or } b \text{ or } c+b \pmod{2}$,\\
$bc = \frac{n + \frac{n}{d}\sqrt{m}}{2} = \frac{n - \frac{n}{d}}{2} + \frac{n}{d} \cdot \frac{1+ \sqrt{m}}{2} \equiv 0 \pmod{2} $.
\end{tabular}
\end{center}
and all the 'or'-s in the above remainders depend only on $m \pmod{8}$ and $d \pmod{4}$, hence we have to check 4 cases.
\\ \\
(a) $m \equiv 1 \pmod{8}$, $d \equiv 1 \pmod{4}$. We construct an ideal
$I = (x^2+x,y^2,z^2,xy+z,xz+z,yz)$ and obtain prime ideals $A = (z,y,x)$ and $B = (z,y,x+1)$ satisfying $A^2 = (z,x)$, $B^2=(y+z,x+1)$, $AB = (y+z,z)$ and $(AB)^2 = 0$. Hence we have $(2) = (\mathfrak{p}_1\mathfrak{p}_2)^2 \implies e=2, \ f=1$.  Because $y \in A, \not\in A^2$ (the same is true for $\mathfrak{p}_2$) we can take $b = \pi \in \mathfrak{p}_1 \setminus \mathfrak{p}_1^2$. Then  $\pi^2 = b^2 = n \equiv 2 \pmod{\mathfrak{p}_1^4}$, hence by Theorem \ref{1.3} we have $s=6$, for $\mathfrak{p}_1$ and $\mathfrak{p}_2$.
\\ \\
(b) $m \equiv 1 \pmod{8}$, $d \equiv 3 \pmod{4}$. We have $I = (x^2+x,y^2,z^2,xy+z+y,xz,yz)$ and two prime ideals $A = (z,y,x)$ and $B=(z,y,x+1)$ with $A^2 = (y+z,x)$, $B^2 = (z,x+1)$, $AB = (y+z,z)$ and $(AB)^2 = 0$. As previously, we obtain $(2) = (\mathfrak{p}_1\mathfrak{p}_2)^2$ thus $e=2, \ f=1$. Take $b = \pi \in \mathfrak{p}_1 - \mathfrak{p}_1^2$ (which is also true for $\mathfrak{p}_2$). Same reasoning as in (a) shows that $s=6$ for both prime ideals.
\\ \\
(c) $m \equiv 5 \pmod{8}$, $d \equiv 1 \pmod{4}$. We have $I = (x^2+x+1,y^2,z^2,xy+z,xz+z+y,yz)$ and a prime ideal $J = (z,y)$ satisfying $J^2=0$. This gives $(2) = \mathfrak{p}^2 \implies e=2,\ f=2.$ Hence, $ s=2$ by Theorem \ref{1.2}.
\\ \\
(d) $m \equiv 5 \pmod{8}$, $d \equiv 3 \pmod{4}$. We have $I = (x^2+x+1,y^2,z^2,xy+z+y,xz+y,yz)$ and a prime ideal $J = (z,y)$ satisfying $J^2=0$. This gives $(2) = \mathfrak{p}^2 \implies e=2,\ f=2$. Hence, $s=2$ by theorem \ref{1.2}.
\\ \\
\textbf{Case 3.4.}  $ m\equiv n \equiv k \equiv 1 \pmod{4}$ and the basis of $\mathcal{O}_K$ is $\{1, \frac{1+\sqrt{m}}{2}, \frac{1+\sqrt{n}}{2},(\frac{1+\sqrt{m}}{2})(\frac{1+\sqrt{k}}{2})  \} = \{1,a,b,c\}$. 
\begin{center}
\begin{tabular}{@{}l@{}}
$a^2 = \frac{m-1}{4} + \frac{1+\sqrt{m}}{2}   $,\\
$b^2 = \frac{n-1}{4} + \frac{1+\sqrt{n}}{2}   $,\\
$c^2  =  \frac{(m-1)(k-1)}{16}  + \frac{m-1}{4} \cdot \frac{m+d}{2d} + \frac{k-m}{4}a + \frac{m}{d} \cdot \frac{1-m}{4}b + \frac{m+1}{2}c  $,\\
$ab = \frac{m-1}{4} + \frac{1-d}{2}a + \frac{1-m}{2}b + dc $,\\
$ac = \frac{m-1}{4} \cdot \frac{m+d}{2d} + \frac{1-m}{4}a + \frac{m}{d} \cdot \frac{1-m}{4}b + \frac{m+1}{2}c  $,\\
$bc = \frac{m-1}{4} \cdot \frac{n+d}{2d} +  \frac{\frac{n}{d}-d}{4}a + \frac{1-m}{4}b+ \frac{d+1}{2}c  $.
\end{tabular}
\end{center}
Because $k = \frac{mn}{d^2} \equiv mn \pmod{8}$ we have to distinguish four cases:
\\ \\
(a) $m \equiv n \equiv k \equiv 1 \pmod{8}$. We have $I = (x^2+x,y^2+y,z^2+z,xy+z,xz+z,yz+z)$. Using Macaulay2 one can find four different prime ideals $A = (z,y,x)$, $B = (z,y,x+1) $, $C = (z,y+1,x)$, $D = (z+1,y+1,x+1) $ such that $ABCD = 0.$ This implies that $ (2) = \mathfrak{p}_1\mathfrak{p}_2\mathfrak{p}_3\mathfrak{p}_4$. Hence for each $\mathfrak{p}_i$, $i=1,2,3,4$ we get $ e=f=1$ and $ s=15$ by Theorem \ref{1.1}.
\\ \\
(b) $m \equiv n \equiv 5,\  k \equiv 1 \pmod{8}$. Modulo 2 we get:
\begin{center}
\begin{tabular}{@{}l@{}}
$a^2 \equiv 1 +a$,\\
$b^2 \equiv 1 + b$,\\
$c^2 \equiv 1 +a + b + c \text{ or }  a + b + c$, \\
$ab \equiv 1+c \text{ or } 1+a+c$, \\
$ac \equiv  1+a + b + c \text{ or }  a + b + c$, \\
$bc \equiv 1+a + b + c \text{ or } a + b $.
\end{tabular}
\end{center}
Where all the 'or'-s depend on $d \pmod{4}$.

If $d \equiv 1 \pmod{4}$ then $I = (x^2+1+x,y^2+1+y,z^2+1+x+y+z,xy+1+z,xz+1+x+y+z,yz+1+x+y+z)$ and using Macaulay2 we get prime ideals $A=(z+1,y)$ and $B = (y+1,x+z+1)$ satisfying $AB = 0$. Hence $  (2) = \mathfrak{p}_1\mathfrak{p}_2 \implies e=1,f=2$ and $s=2$ by Theorem \ref{1.2}.

If $d \equiv 3 \pmod{4}$ then $I = (x^2+1+x,y^2+1+y,z^2+x+y+z,xy+1+x+z,xz+x+y+z,yz+x+y) $ and we get prime ideals $A = (z,x+y)$ and $B = (y+z+1,x+z)$ such that $AB = 0$. This implies $ (2) = \mathfrak{p}_1\mathfrak{p}_2 $ hence, $ e=1,f=2$ and $ s=2$ by Theorem \ref{1.2}.
\\ \\
The last remaining cases: (c) $m \equiv k \equiv 5, n \equiv 1 \pmod{8}$ and (d) $n \equiv k \equiv 5, m \equiv 1 \pmod{8}$ follow directly from rearranging $n,m,k$. In both cases we have $s=2$.

\end{proof}

\end{document}
